% ID: SEL-D0038-OA-E02
% ESTADO: validada. S039, cotejada con examen y criterios.
\textbf{a)} En $-1$, el primer tramo da $f(-1)=a/2$,
que coincide con el límite izquierdo. El límite derecho
es $-1-3-6-2=-12$. Igualar ambos da $\boxed{a=-24}$,
valor para el que el límite coincide con $f(-1)$.

\textbf{b)} La condición de pertenencia es $g(1)=2$, que da
$b+c=3$. Como $g''(x)=6x+2b$, la inflexión exige $6+2b=0$,
de donde $b=-3$ y $c=6$. Con estos valores,
$g''(x)=6(x-1)$ cambia de signo en $1$, así que la inflexión
queda comprobada. Resulta $\boxed{b=-3,\ c=6}$.
